\section{Adaptive Tree Objective Proofs}
\label{app:adaptive-tree-proofs}

This appendix isolates the adaptive-tree optimizer-transfer results used in the
main text. The point of separating them is organizational: they rely on the
same basic preference-learning machinery as the earlier DPO and GRPO sections,
but they answer a different question. The earlier appendices show that
summarization preserves or approximately preserves a given objective. This
appendix shows that one may \emph{optimize} a stochastic tree objective
directly, even when the oracle is only an approximation to the true latent
target.

\subsection{Expected Tree Objectives and Oracle Uncertainty}

Training under adaptive tree construction often targets the \emph{expected tree
objective}
\[
\widehat{\mathcal{L}}_{\tau}(\theta)
:=
\E_{T \sim \tau(x)}[\mathcal{L}_{\mathrm{tree}}(\theta;T)],
\]
where $\tau(x)$ is a distribution over admissible trees for document $x$.

\begin{theorem}[Expected-tree optimizer transfer]
\label{thm:app-expected-tree-transfer}
Let $\Theta_{\mathrm{adm}}$ be any admissible parameter class. Suppose that for
each $\theta \in \Theta_{\mathrm{adm}}$,
\[
\E_{T \sim \tau(x)}
\left[
  \left|
    \mathcal{L}_{\mathrm{true}}(\theta) -
    \mathcal{L}_{\mathrm{tree}}(\theta;T)
  \right|
\right]
\le \varepsilon_{\tau}(\theta).
\]
If $\widehat{\theta} \in \arg\min_{\theta \in \Theta_{\mathrm{adm}}}
\widehat{\mathcal{L}}_{\tau}(\theta)$, then
\[
\mathcal{L}_{\mathrm{true}}(\widehat{\theta})
\le
\mathcal{L}_{\mathrm{true}}(\theta)
  + \varepsilon_{\tau}(\widehat{\theta})
  + \varepsilon_{\tau}(\theta)
\quad
\text{for all } \theta \in \Theta_{\mathrm{adm}}.
\]
\end{theorem}

\begin{proof}
Write
$\widehat{\mathcal{L}}_{\tau}(\theta)
= \E_{T}[\mathcal{L}_{\mathrm{tree}}(\theta;T)]$.
By Jensen,
\[
\left|
\mathcal{L}_{\mathrm{true}}(\theta) - \widehat{\mathcal{L}}_{\tau}(\theta)
\right|
\le
\E_T\!\left[
\left|
\mathcal{L}_{\mathrm{true}}(\theta)-\mathcal{L}_{\mathrm{tree}}(\theta;T)
\right|
\right]
\le \varepsilon_{\tau}(\theta).
\]
So the expected tree objective is a nonuniform perturbation of the true
objective. Applying the standard perturbation argument to an exact minimizer of
$\widehat{\mathcal{L}}_{\tau}$ gives
\[
\widehat{\mathcal{L}}_{\tau}(\widehat{\theta})
\le
\widehat{\mathcal{L}}_{\tau}(\theta)
\]
for every admissible $\theta$, and the two one-sided perturbation bounds then
yield
\[
\mathcal{L}_{\mathrm{true}}(\widehat{\theta})
\le
\widehat{\mathcal{L}}_{\tau}(\widehat{\theta})
  + \varepsilon_{\tau}(\widehat{\theta})
\le
\widehat{\mathcal{L}}_{\tau}(\theta)
  + \varepsilon_{\tau}(\widehat{\theta})
\le
\mathcal{L}_{\mathrm{true}}(\theta)
  + \varepsilon_{\tau}(\widehat{\theta})
  + \varepsilon_{\tau}(\theta),
\]
which is the desired claim.
\end{proof}

\paragraph{Where the slack comes from.}
In our setting the quantity $\varepsilon_{\tau}(\theta)$ decomposes naturally as
\[
\varepsilon_{\tau}(\theta)
=
\E_T[\varepsilon_{\mathrm{oracle}}(\theta,T)]
\;+\;
\E_T[\varepsilon_{\mathrm{transport}}(T)].
\]
The first term measures oracle misspecification relative to the true latent
target. The second measures objective drift caused by summarization and tree
reduction. The total optimization slack is therefore exactly
``oracle misspecification plus summary-induced drift,'' averaged over the tree
policy.

\paragraph{Method-specific corollaries.}
The abstract theorem specializes directly to the three preference-learning
families studied in the paper:
\begin{itemize}[nosep]
    \item \textbf{DPO}: the transport term is
    $2|\beta|L_{\mathrm{pol}}$ times the expected local-law budget.
    \item \textbf{GRPO-PL}: the transport term is
    $L_{\mathrm{grpo}}$ times the expected local-law budget.
    \item \textbf{GRPO-RL}: the transport term is
    $L_{\mathrm{grpo-rl}}$ times the expected local-law budget.
\end{itemize}

\begin{theorem}[High-probability expected-tree optimizer transfer]
\label{thm:app-expected-tree-transfer-hp}
Let $G$ be a good event with $\Pr(G^c)\le \delta$, and let
$\widehat{\theta}(\omega)$ be a data-dependent optimizer selection rule. Assume
that whenever $\omega \in G$, the selected parameter
$\widehat{\theta}(\omega)$ exactly minimizes the expected tree objective
$\widehat{\mathcal{L}}_{\tau}$ over the admissible class. If
\[
\E_{T \sim \tau(x)}
\left[
  \left|
    \mathcal{L}_{\mathrm{true}}(\theta) -
    \mathcal{L}_{\mathrm{tree}}(\theta;T)
  \right|
\right]
\le \varepsilon_{\tau}(\theta)
\]
for every admissible $\theta$, then
\[
\Pr\!\left(
\widehat{\theta}(\omega)
\notin
\varepsilon_{\tau}\text{-}\arg\min \mathcal{L}_{\mathrm{true}}
\right)
\,\le\,
\delta.
\]
Equivalently, with probability at least $1-\delta$, the selected optimizer
satisfies
\[
\mathcal{L}_{\mathrm{true}}(\widehat{\theta}(\omega))
\le
\mathcal{L}_{\mathrm{true}}(\theta)
  + \varepsilon_{\tau}(\widehat{\theta}(\omega))
  + \varepsilon_{\tau}(\theta)
\quad
\text{for all admissible } \theta.
\]
\end{theorem}

\begin{proof}
On the event $G$, Theorem~\ref{thm:app-expected-tree-transfer} applies
pointwise to the selected optimizer $\widehat{\theta}(\omega)$, so failure of
the displayed near-optimality statement is impossible on $G$. Therefore the
failure event is contained in $G^c$, and its probability is at most
$\Pr(G^c)\le \delta$.
\end{proof}

\paragraph{Interpretation.}
This is the form used with audits, held-out certification events, or any other
data-dependent procedure that certifies the local-law or transport budget only
with high probability. Confidence uncertainty contributes only through the
failure probability of the certifying event, while optimization slack still
decomposes additively into expected oracle uncertainty and expected
tree-transport error.
